%
% Copyright (c) 2015 Radoslaw Kujawa.
% All rights reserved.
%
% Redistribution and use in source and binary forms, with or without
% modification, are permitted provided that the following conditions
% are met:
%
% 1. Redistributions of source code must retain the above copyright
%    notice, this list of conditions and the following disclaimer.
% 2. Redistributions in binary form must reproduce the above copyright
%    notice, this list of conditions and the following disclaimer in the
%    documentation and/or other materials provided with the distribution.
%
% THIS SOFTWARE IS PROVIDED BY RADOSLAW KUJAWA (THE AUTHOR) AND CONTRIBUTORS
% ``AS IS'' AND ANY EXPRESS OR IMPLIED WARRANTIES, INCLUDING, BUT NOT LIMITED
% TO, THE IMPLIED WARRANTIES OF MERCHANTABILITY AND FITNESS FOR A PARTICULAR
% PURPOSE ARE DISCLAIMED.  IN NO EVENT SHALL THE AUTHOR OR CONTRIBUTORS
% BE LIABLE FOR ANY DIRECT, INDIRECT, INCIDENTAL, SPECIAL, EXEMPLARY, OR
% CONSEQUENTIAL DAMAGES (INCLUDING, BUT NOT LIMITED TO, PROCUREMENT OF
% SUBSTITUTE GOODS OR SERVICES; LOSS OF USE, DATA, OR PROFITS; OR BUSINESS
% INTERRUPTION) HOWEVER CAUSED AND ON ANY THEORY OF LIABILITY, WHETHER IN
% CONTRACT, STRICT LIABILITY, OR TORT (INCLUDING NEGLIGENCE OR OTHERWISE)
% ARISING IN ANY WAY OUT OF THE USE OF THIS SOFTWARE, EVEN IF ADVISED OF THE
% POSSIBILITY OF SUCH DAMAGE.
%
% 
\documentclass[dvipsnames,table]{beamer}
\usepackage{polski}

\usetheme{Rochester}
\usecolortheme{orchid}

\usepackage{listings}
\usepackage{ucs}
\usepackage[utf8x]{inputenc}
\usepackage{wasysym}
\usepackage[normalem]{ulem}
\usepackage{amsmath}
\usepackage{hyperref}

\setbeamertemplate{navigation symbols}{}
\setbeamertemplate{caption}[numbered]
\setbeamerfont{caption}{size=\scriptsize}
\setbeamercolor{framenote}{bg=OSEC-red!25}
\setbeamercolor{rednote}{bg=Red!25}
\setbeamercolor{palette primary}{use=structure,fg=white,bg=OSEC-red}
\setbeamercolor{palette secondary}{use=structure,fg=white,bg=OSEC-red2}

\setbeamertemplate{itemize item}{\scriptsize\raise1pt\hbox{\donotcoloroutermaths$\blacktriangleright$}}
\setbeamertemplate{itemize subitem}{\tiny\raise1pt\hbox{\donotcoloroutermaths$\bullet$}}
\setbeamertemplate{itemize subsubitem}{\tiny\raise1pt\hbox{\donotcoloroutermaths{--}}}

\setbeamertemplate{enumerate item}{\insertenumlabel.}
\setbeamertemplate{enumerate subitem}{\insertenumlabel.\insertsubenumlabel}
\setbeamertemplate{enumerate subsubitem}{\insertenumlabel.\insertsubenumlabel.\insertsubsubenumlabel}
\setbeamertemplate{enumerate mini template}{\insertenumlabel}

\setbeamercolor{itemize item}{fg=OSEC-red, bg=OSEC-red}
\setbeamercolor{itemize subitem}{fg=OSEC-red, bg=OSEC-red}
\setbeamercolor{itemize subsubitem}{fg=OSEC-red, bg=OSEC-red}

\setbeamercolor{section number projected}{fg=white,bg=OSEC-red}
\setbeamercolor{subsection number projected}{fg=white,bg=OSEC-red}
\setbeamercolor{button}{bg=OSEC-red,fg=white}

\setbeamertemplate{section in toc}[circle]
\setbeamertemplate{subsection in toc}[square]


\definecolor{OSEC-red}{RGB}{160,29,44}
\definecolor{OSEC-red2}{RGB}{177,76,12}
\hypersetup{colorlinks=true,linkcolor=white,urlcolor=OSEC-red}

\setlength{\tabcolsep}{8pt}
\renewcommand{\arraystretch}{1.2}

\newcommand{\nbsdcolor}[1] {
	{\color{OSEC-red} #1}
}

\lstset{
   language=java,
   basicstyle=\tiny,
   breaklines=true,
   escapechar=\@,
   commentstyle=\color{OSEC-red}
}

\AtBeginSection[]{
\frame{

\begin{center}

{\usebeamerfont{section name}\usebeamercolor[fg]{section name}Część~\insertsectionnumber}
    \vskip1em\par

	\begin{beamercolorbox}[sep=12pt,center]{palette primary}
		\usebeamerfont{section title}\insertsection\par
	\end{beamercolorbox}
\end{center}

%\sectionpage
}
}

\title{Wykorzystanie rekordów DNS \\ do weryfikacji kluczy publicznych SSH}


\author{Radosław Kujawa -- radoslaw.kujawa@osec.pl}

\institute{OSEC}

\begin{document}

\begin{frame}
\titlepage
\end{frame}

\begin{frame}[allowframebreaks]
\frametitle{Spis treści}
{
\hypersetup{colorlinks=true,linkcolor=black,urlcolor=OSEC-red}
\tableofcontents
}
\end{frame}


\section{Zagrożenia bezpieczeństwa związane z weryfikacją kluczy publicznych hostów}

\begin{frame}
\frametitle{Pierwsze połączenie}
\begin{itemize}
	\item Każdy z nas używa SSH...
	\item Ale czy każdy z nas zdaję sobie sprawę z zagrożeń? 
\end{itemize}
\end{frame}

\section{Jak wykorzystać rekordy SSHFP do zautomatyzowanej weryfikacji kluczy}

\begin{frame}
\frametitle{Problem - weryfikacja kluczy}
\begin{itemize}
	\item Foo {\tt bar}
	\item \href{http://example.com/}{example}
\end{itemize}
\end{frame}

\section{Konfiguracja klienta SSH}

\begin{frame}
\frametitle{Konfiguracja klienta SSH}
\begin{itemize}
	\item Statystyki pracy serwera dostępne przez API
\end{itemize}
\end{frame}

\begin{frame}
\frametitle{Cechy nowego API}
\begin{itemize}
	\item Prostota i funkcjonalność
	\item Oparte o komunikaty JSON
	\item ,,Zdetypowane'' -- zmiany możliwe z innej maszyny nawet bez natywnych bibliotek serwera, każdy klient HTTP może służyć do zarządznia serwerem
	\item Pozwala na przesyłanie nie tylko konfiguracji, ale także artefaktów na zdalne maszyny
\end{itemize}
\end{frame}

\begin{frame}[fragile]
\frametitle{Foo foo foo}
\begin{verbatim}
bar bar bar
\end{verbatim}
\end{frame}

\section{Automatyzacja procesu aktualizacji stref}

\section{Zabezpieczanie DNS przed spoofingiem}

\begin{frame}
\frametitle{Koniec\ldots}
%\vspace*{0.15cm}
\begin{center}
\includegraphics[scale=0.5]{img-oseclogo.png}

Dziękuje!

Czy są pytania?
\end{center}

\end{frame}

 
\end{document}
